% The Monte Carlo uncertainty method, written out.
% Build with:  cd docs && pdflatex monte_carlo_method.tex && pdflatex monte_carlo_method.tex
%
% Companion to notebooks/monte_carlo_explained.ipynb, which runs the same method
% line by line. This document carries the mathematics; the notebook carries the
% executable version and verifies it reproduces the pipeline exactly.
\documentclass[11pt,a4paper]{article}

\usepackage[margin=1in]{geometry}
\usepackage{amsmath,amssymb}
\usepackage{booktabs}
\usepackage{array}
\usepackage{graphicx}
\usepackage{xcolor}
\usepackage[T1]{fontenc}
\usepackage{lmodern}
\usepackage[colorlinks=true,linkcolor=linknavy,urlcolor=linknavy]{hyperref}
\emergencystretch=3em

\definecolor{linknavy}{HTML}{2A78D6}
\definecolor{inkSoft}{HTML}{52514E}
\definecolor{boxbg}{HTML}{F4F6FA}
\definecolor{boxrule}{HTML}{C9D6EA}
\definecolor{warnbg}{HTML}{FDF3E7}
\definecolor{warnrule}{HTML}{E3B778}

\newcommand{\keybox}[2]{%
  \par\vspace{6pt}\noindent
  \fcolorbox{boxrule}{boxbg}{%
    \begin{minipage}{0.965\linewidth}\vspace{4pt}
    {\bfseries #1}\par\vspace{3pt}#2\vspace{5pt}
    \end{minipage}}\par\vspace{6pt}}

\newcommand{\warnbox}[2]{%
  \par\vspace{6pt}\noindent
  \fcolorbox{warnrule}{warnbg}{%
    \begin{minipage}{0.965\linewidth}\vspace{4pt}
    {\bfseries #1}\par\vspace{3pt}#2\vspace{5pt}
    \end{minipage}}\par\vspace{6pt}}

\title{\vspace{-1.4cm}Monte Carlo uncertainty on a predicted\\ lattice thermal conductivity}
\author{The method used in this project, written out}
\date{}

\begin{document}
\maketitle
\vspace{-0.8cm}

\noindent
\textit{Companion to} \texttt{notebooks/monte\_carlo\_explained.ipynb}, \textit{which
implements every equation below and verifies it reproduces}
\texttt{scripts/cgcnn/07\_predict\_kappa.py} \textit{exactly.}

% =====================================================================
\section{The problem}

The pipeline predicts two elastic moduli and turns them into a lattice thermal
conductivity through a chain of closed-form relations:

\begin{align}
v_l &= \sqrt{\frac{K + \tfrac{4}{3}G}{\rho}}, \qquad
v_t = \sqrt{\frac{G}{\rho}}, \qquad
v_s = \left[\frac{1}{3}\left(\frac{1}{v_l^{3}} + \frac{2}{v_t^{3}}\right)\right]^{-1/3} \\[4pt]
\nu &= \frac{(v_l/v_t)^2 - 2}{2(v_l/v_t)^2 - 2}, \qquad
\gamma = \frac{3(1+\nu)}{2(2-3\nu)} \\[4pt]
\kappa_L &= \frac{G\,v_s\,V^{1/3}}{N\,T}\,e^{-\gamma}
\end{align}

Every one of these is exact. But $K$ and $G$ are \emph{predictions}, and a prediction has
an uncertainty. The question is what that uncertainty does to $\kappa_L$.

\subsection{Why the textbook answer does not work here}

The standard tool is the \emph{delta method}: linearise about the prediction and combine
partial derivatives,
\[
\sigma_{\kappa}^2 \;\approx\;
\left(\frac{\partial \kappa_L}{\partial K}\right)^{\!2}\sigma_K^2 +
\left(\frac{\partial \kappa_L}{\partial G}\right)^{\!2}\sigma_G^2 .
\]
This is a first-order Taylor expansion, and it is only as good as the assumption that
$\kappa_L$ is close to linear over the width of the uncertainty. Ours is not. Reading the
chain from the inside out: $K$ and $G$ enter $v_l/v_t$ as a ratio under square roots; that
ratio enters $\nu$ through a rational function; $\nu$ enters $\gamma$ through another
rational function with $(2-3\nu)$ in the denominator; and $\gamma$ then enters
\emph{exponentially}. Composing four non-linear steps and expecting a linear response over
a $\pm 15\%$ input range is not reasonable, and the derivatives themselves are unpleasant
to write down correctly.

\keybox{The Monte Carlo alternative}{Instead of differentiating the function, evaluate it
many times on inputs drawn to represent the uncertainty, and look at the spread of the
outputs. No derivatives, no linearity assumption, and the answer becomes exact as the
number of draws grows. The cost is arithmetic, which is cheap.}

% =====================================================================
\section{What the uncertainty is, and where it comes from}

The moduli come from an \textbf{ensemble}: three networks trained from different random
initialisations on the same data and split. For crystal $i$ they produce three values,
whose mean is the prediction and whose standard deviation is the uncertainty. Both are
computed on $\log_{10}$ of the modulus, because that is the space the networks are trained
in:
\[
\mu_{K,i} = \log_{10} \hat{K}_i,
\qquad
\sigma_{K,i} = \operatorname{sd}\big(\log_{10} K_i^{(1)}, \log_{10} K_i^{(2)}, \log_{10} K_i^{(3)}\big),
\]
and likewise for $G$. In this project's 1{,}213-crystal prediction set the median
$\sigma_K$ is $0.0235$ and the median $\sigma_G$ is $0.0295$, i.e.\ roughly $5$--$7\%$
multiplicative uncertainty.

% =====================================================================
\section{The method}

\subsection{Step 1 --- draw}

For each crystal $i$ and each sample $j = 1 \dots M$,
\begin{align}
\log_{10} K_i^{(j)} &= \mu_{K,i} + \sigma_{K,i}\, z^{(j)}_{K,i},
&& z_{K} \sim \mathcal{N}(0,1) \\
\log_{10} G_i^{(j)} &= \mu_{G,i} + \sigma_{G,i}\, z^{(j)}_{G,i},
&& z_{G} \sim \mathcal{N}(0,1)
\end{align}
then $K_i^{(j)} = 10^{\log_{10} K_i^{(j)}}$ and similarly for $G$. This is the
location--scale construction: a standard normal scaled by $\sigma$ and shifted by $\mu$ is
distributed $\mathcal{N}(\mu, \sigma^2)$.

Two modelling choices are embedded here and both deserve to be stated rather than assumed.

\medskip
\noindent\textbf{Why a normal distribution.} Three members give a mean and a spread and
essentially no information about shape. Among all distributions with a given mean and
variance, the normal is the maximum-entropy choice --- the one that adds the least
unwarranted structure.

\medskip
\noindent\textbf{Why $z_K$ and $z_G$ are independent.} The $K$ ensemble and the $G$
ensemble are separately seeded runs sharing no weights and no random state, so nothing
links their draws. \emph{This is the weakest assumption in the method.} Both ensembles see
the same crystal, and a crystal that is hard for one is plausibly hard for the other; a
positive correlation between the two errors would widen the true interval beyond what this
produces.

\subsection{Step 2 --- propagate}

Every sampled pair goes through the \emph{same} function used for the point estimate:
\[
\kappa_{L,i}^{(j)} = f\!\left(K_i^{(j)},\, G_i^{(j)};\; V_i, \rho_i, \bar{M}_i, N_i\right),
\]
where $f$ is equations (1)--(3) and the structural quantities are held \textbf{fixed}: they
are read from the CIF, not predicted, so they carry no uncertainty of this kind.

Because the point estimate and the interval come from the same $f$, they are consistent by
construction rather than by a separate approximation --- the median of the samples cannot
drift away from the reported central value for reasons of method.

\subsection{Step 3 --- summarise}

With $M$ values of $\kappa_L$ per crystal, report the empirical quantiles
\[
p_{05,i} = Q_{0.05}\big(\{\kappa_{L,i}^{(j)}\}\big), \quad
p_{50,i} = Q_{0.50}\big(\cdot\big), \quad
p_{95,i} = Q_{0.95}\big(\cdot\big).
\]

\keybox{Why quantiles, and not mean $\pm 2\sigma$}{The input distribution is symmetric; the
output is not. On the worked example in the notebook (LaGa) the median sits at
$4.48$~W\,m$^{-1}$K$^{-1}$, with $2.31$ down to $p_{05}$ and $4.68$ up to $p_{95}$ --- the
upper half of the interval is twice the lower. A symmetric interval would misstate the
answer in both directions. Quantiles assume nothing about shape.}

% =====================================================================
\section{Why everything happens in $\log_{10}$}

This is the design choice that is easiest to overlook and most consequential.

\medskip
\noindent\textbf{The error is multiplicative.} The networks are trained on
$\log_{10}(\text{modulus})$ against a log-space loss, so their error is naturally a
percentage, not a number of GPa. An error of $12$~GPa is catastrophic on a $20$~GPa crystal
and negligible on a $400$~GPa one; an error of $15\%$ means the same thing on both.
Sampling in log space encodes that correctly.

\medskip
\noindent\textbf{Negative moduli become impossible.} Since $10^x > 0$ for every real $x$, a
log-space draw can never be non-positive. Sampling additively in GPa can, and does: on this
project's softest crystal (Rb, $\hat{K} = 2.01$~GPa, $\sigma_K = 0.1285$), an
equivalent-width additive draw of 20{,}000 samples produced 5 negative moduli. A negative
modulus is unphysical, and it makes $\sqrt{G/\rho}$ undefined --- the Monte Carlo would
fail on exactly the soft crystals the screen exists to find.

% =====================================================================
\section{How many samples}

Monte Carlo error falls as $M^{-1/2}$, so accuracy is cheap but improves slowly. This
project uses $M = 2000$. Measured convergence on the worked crystal, against a
20{,}000-sample reference:

\begin{center}
\begin{tabular}{@{}rrr@{}}
\toprule
$M$ & $p_{05}$ & $p_{95}$ \\
\midrule
50      & 2.466 & 8.517 \\
250     & 2.039 & 9.199 \\
1{,}000 & 2.234 & 9.316 \\
\textbf{2{,}000} & \textbf{2.172} & \textbf{9.160} \\
5{,}000 & 2.155 & 9.356 \\
20{,}000 & 2.169 & 9.339 \\
\bottomrule
\end{tabular}
\end{center}

\noindent
At $M = 2000$, $p_{05}$ is within $0.13\%$ and $p_{95}$ within $1.92\%$ of the reference.
Both are far smaller than the interval itself, which spans a factor of $4.2$. Increasing
$M$ further would refine a number whose real uncertainty is dominated by something else
entirely --- which is the subject of the next section.

% =====================================================================
\section{The interval is too narrow, and by how much}

Everything above is arithmetically correct and still produces an interval that is badly
overconfident. This is the most important result in this document.

Step 25 of the pipeline tested it on the held-out set, where the true DFT moduli are known,
and asked: of the crystals where the method claimed $X\%$ confidence, what fraction did the
interval actually contain?

\begin{center}
\begin{tabular}{@{}rrrrr@{}}
\toprule
nominal & \multicolumn{2}{c}{$K_{\mathrm{VRH}}$} & \multicolumn{2}{c}{$G_{\mathrm{VRH}}$} \\
\cmidrule(lr){2-3}\cmidrule(lr){4-5}
 & raw & recalibrated & raw & recalibrated \\
\midrule
0.5 & 0.234 & 0.496 & 0.214 & 0.512 \\
0.8 & 0.414 & 0.789 & 0.394 & 0.811 \\
\textbf{0.9} & \textbf{0.508} & 0.895 & \textbf{0.479} & 0.909 \\
\bottomrule
\end{tabular}
\end{center}

\noindent
A nominal $90\%$ interval covers about $\mathbf{50\%}$ of true values.

\subsection{Why}

The Monte Carlo faithfully propagates the ensemble's \emph{disagreement}. But disagreement
is only the variance component of the error:
\[
\mathbb{E}\big[(\hat{y}-y)^2\big] \;=\; \underbrace{\big(\mathbb{E}[\hat{y}]-y\big)^2}_{\text{bias}^2} \;+\; \underbrace{\operatorname{Var}(\hat{y})}_{\text{what the ensemble sees}} .
\]
Three networks with the same architecture, the same training data and the same split share
their blind spots. When all three are wrong in the same direction, the bias term is large
and the variance term is small --- the ensemble agrees with itself, the sampled spread is
narrow, and the method reports high confidence in a wrong answer. Measured here, member
disagreement runs $4$--$6\times$ smaller than the true error.

\warnbox{The generalisable lesson}{\textbf{An ensemble's spread is a lower bound on its
uncertainty, never the whole of it.} It must be calibrated against held-out truth before it
is quoted as a confidence interval. Rescaling $\sigma$ by a factor fitted on validation data
restores $\sim$90\% coverage out of sample (right-hand columns above) --- but the factor has
to be measured, and it is level-dependent: roughly $4.2\times$ at nominal $0.8$ and
$7.0\times$ at $0.9$. Applying that correction to this project's DFT shortlist collapsed
almost every ``tight'' confidence label to ``wide''.}

% =====================================================================
\section{Summary}

\begin{center}
\begin{tabular}{@{}p{0.13\linewidth}p{0.40\linewidth}p{0.40\linewidth}@{}}
\toprule
step & what happens & why \\
\midrule
draw & $M = 2000$ pairs from $\mathcal{N}(\mu, \sigma^2)$ in $\log_{10}$ space & the error is multiplicative, and $10^x$ cannot be negative \\[4pt]
propagate & every pair through the same $f$ as the point estimate & no linearisation, no derivatives; interval consistent with the central value by construction \\[4pt]
summarise & the $5$th, $50$th and $95$th percentiles & the output is skewed although the input was symmetric \\
\bottomrule
\end{tabular}
\end{center}

\noindent
And the caveat that outweighs the rest: this propagates \emph{variance} only. Uncorrected,
the interval covers about half of what it claims.

\end{document}
